% !TeX root = LiteraturosApzvalga_2.tex


\section{K artimiausių kaimynų metodas}
%==================================================================

Šiame skyriuje pristatomas KNN metodas.
Apibūdinamas veikimo principas, pagrindiniai metodą apibrėžiantys sprendimai, tokie kaip atstumo funkcija, kaimynų skaičius ir prognozės sudarymas, atsisakymo mechanizmas.

%------------------------------------------------------------------

\subsection{Veikimo principas}
%------------------------------------------------------------------

Metodas grindžiamas prielaida, kad panašiems pavyzdžiams būdingos panašios baigtys.
Užuot mokęsis apibendrintos priklausomybės tarp įvesties ir išvesties, metodas prognozuojamą objektą tiesiogiai lygina su istoriniais pavyzdžiais.
Veikimą galima aprašyti keturiais žingsniais:
\begin{enumerate}
    \item Vertinamas objektas aprašomas požymių vektoriumi.
    \item Apskaičiuojamas atstumas nuo šio vektoriaus iki visų mokymo pavyzdžių.
    \item Mokymo pavyzdžiai išrikiuojami pagal atstumą ir atrenkami $k$ artimiausių -- kaimynai.
    \item Iš kaimynų reikšmių sudaroma prognozė.
\end{enumerate}
Klasifikuojant kaimynai balsuoja už klasę, o regresijoje jų reikšmės vidurkinamos.
Metodo veikimas dvimatėje požymių erdvėje pavaizduotas \ref{fig:knn-kaimynai} pav.
Laiko eilučių prognozėje istoriniai pavyzdžiai yra praeities langai, o prognozė sudaroma iš reikšmių, einančių po panašių istorinių langų.

\input{paveikslas_knn_kaimynai}

Artimiausių kaimynų taisyklę klasifikavimo uždaviniui dar 1951 metais pasiūlė Fix ir Hodges \cite{fix1951}, o teorinį pagrindą jai suteikė Cover ir Hart \cite{cover1967}.
Jie įrodė, kad imties dydžiui neribotai augant, vieno artimiausio kaimyno taisyklės klaidos tikimybė $R$ neviršija dvigubos mažiausios teoriškai įmanomos Bajeso klaidos $R^*$, tai yra $R^* \le R \le 2R^*(1-R^*)$.
Vadinasi, joks kitas metodas, turėdamas tuos pačius duomenis, klaidos neįstengtų sumažinti daugiau nei perpus.
Todėl autoriai teigia, kad vieno artimiausio kaimyno metode slypi pusė visos begalinėje imtyje esančios klasifikavimo informacijos \cite{cover1967}.
Nors teorinis įrodymas gautas klasifikavimo uždaviniams, pati KNN idėja praktikoje plačiai taikoma ir regresijai bei laiko eilučių prognozavimui.

Skirtingai nuo daugumos mašininio mokymosi metodų, KNN neturi mokymo etapo.
\cite{wettschereck1997} tokius metodus vadina tingiuoju mokymusi (angl. \textit{lazy learning}). Metodas tiesiog saugo visą mokymo aibę, o pronozės skaičiavimą atideda iki užklausos pateikimo momento.
KNN yra neparametrinis metodas, reiškiantis, kad jis nedaro išankstinių prielaidų apie duomenų pasiskirstymą.
Taip pat, kadangi algoritmas neturi permokymo etapo, nauji duomenys panaudojami iš karto, o kiekvieną prognozė yra interpretuojama, nes ją galima paaiškinti parodant kaimynus, iš kurių ji sudaryta.
Tačiau šis metodas turi ir trūkumų.
Kiekvienai prognozei atstumas skaičiuojamas iki visų istorinių pavyzdžių, todėl skaičiavimo kaštai tenka prognozės momentui ir auga kartu su istorijos ilgiu \cite{halder2024}.
Didelės dimensijos duomenyse metodo efektyvumas mažėja, nes atstumai tarp taškų praranda skiriamąją gebą \cite{halder2024}.
Taip pat prognozė yra istorinių reikšmių vidurkis, todėl metodas negali prognozuoti už istorinių stebėjimų ribų išeinančių reikšmių \cite{martinez2019}.

Dauguma KNN tobulinimo darbų nagrinėja paieškos efektyvumą.
KNN metodų apžvalgoje \cite{halder2024} apžvelgiamas 31 kaimynų paieškos ir 12 jungimo metodų didelės apimties duomenims -- indeksavimo struktūros, atminties ir įvesties--išvesties optimizacijos, lygiagretūs ir paskirstyti skaičiavimai.
Šie sprendimai gerina metodo greitį, bet ne prognozavimo tikslumą.
Šioje literatūros apžvalgoje dėmesys skiriamas prognozavimo mechanizmo tobulinimui: atstumo funkcijų įtakai, kaimynų skaičiui, kaimynų ir požymių svoriams.

%------------------------------------------------------------------

\subsection{Atstumo funkcijos}
%------------------------------------------------------------------

Atstumo funkcija apibrėžia, ką metodas laiko panašumu, todėl nuo jos tiesiogiai priklauso, kurie istoriniai pavyzdžiai bus atrinkti kaimynais.
Dažniausiai naudojamas Euklido atstumas:

\begin{equation}
    d(x, u) = \sqrt{\sum_{j=1}^{m} (x_j - u_j)^2}
    \label{eq:euklidas}
\end{equation}

čia $x = (x_1, \dots, x_m)$ -- dabartinę būseną aprašantis požymių vektorius, $u = (u_1, \dots, u_m)$ -- istorinis požymių vektorius, $m$ -- požymių skaičius.

Prieš skaičiuojant atstumą požymiai normalizuojami, pavyzdžiui į intervalą $[0, 1]$ arba $[-1, 1]$.
Be normalizavimo didelio mastelio požymiai, pavyzdžiui azoto oksidų koncentracijos, nustelbtų mažo mastelio požymius, tokius kaip vėjo greitis, vien dėl matavimo vienetų.
Vienmatėje laiko eilutėje, kai visą požymių vektorių sudaro to paties kintamojo reikšmės, visi požymiai ir taip yra to paties mastelio, todėl normalizavimas nebūtinas \cite{martinez2019}.
Daugiamačiu atveju normalizavimas dažniausiai yra neišvengiamas.

Euklido atstumo funkcijoje visi požymiai atstumui daro vienodą įtaką.
Svertinis Euklido atstumas šios prielaidos atsisako -- kiekvienam požymiui priskiriamas svoris $w_j \geq 0$:

\begin{equation}
    d_w(x, u) = \sqrt{\sum_{j=1}^{m} w_j (x_j - u_j)^2}
    \label{eq:svertinis-euklidas}
\end{equation}

Svoriai leidžia skirtingiems požymiams daryti skirtingą įtaką kaimynų atrankai.
Nereikšmingiems ar triukšmingiems požymiams gali būti priskiriami maži svoriai, o reikšmingiems dideli.
Nulinis svoris požymį iš atstumo pašalina, todėl požymių atranka yra atskiras svertinio atstumo atvejis.

Vis dėlto atskiri svoriai kiekvieną požymį vertina atskirai, nors meteorologiniai požymiai dažnai veikia kartu, pavyzdžiui vėjo greičio įtaka taršos sklaidai
priklauso nuo vėjo krypties.
Bendresnis svertinio atstumo variantas yra Mahalanobio atstumas:

\begin{equation}
    d_M(x, u) = \sqrt{\sum_{j=1}^{m} \sum_{k=1}^{m} M_{jk} (x_j - u_j)(x_k - u_k)}
    \label{eq:mahalanobis}
\end{equation}

čia $M_{jk}$ -- matricos $M$ elementas $j$-oje eilutėje ir $k$-ame stulpelyje.
Įstrižainės elementai $M_{jj}$ atitinka svorius $w_j$ iš (\ref{eq:svertinis-euklidas}) formulės, o
neįstrižaininiai elementai $M_{jk}$, kai $j \neq k$, įvertina požymių $j$ ir $k$ tarpusavio sąveiką.

Taigi vietoje atskirų svorių naudojama pilna svorių matrica $M$, įvertinanti ir visų požymių tarpusavio sąveikas.
Atstumo metrikos mokymosi metodai daugiausia sukurti klasifikavimui \cite{weinberger2009}, tačiau egzistuoja ir regresijai skirtų variantų \cite{weinberger2007mlkr}.
Pilna matrica turi daugiausiai galios, tačiau turi $m^2$ parametrų, tad svorių skaičius auga kvadratiškai nuo požymių skaičiaus, o atskiri matricos elementai nebeturi vieno požymio svarbos prasmės.

Laiko eilutėms sukurta ir specializuotų, elastinių atstumo matų, pavyzdžiui, dinaminis laiko kraipymas (angl. \textit{dynamic time warping}, DTW), kurie leidžia lyginti laike pasislinkusias ar ištemptas sekas.
Laiko eilučių atstumo matų palyginime \cite{paparrizos2020} 71 matas vertintas 128 duomenų rinkiniuose pagal vieno artimiausio kaimyno klasifikavimo tikslumą.
Palyginimas parodė, kad taškas su tašku lyginantys matai, tarp jų Euklido atstumas, yra sparčiausi, tačiau pagal tikslumą kai kurie iš jų, ypač Manheteno atstumo, Euklido atstumą lenkia, o dauguma elastinių matų statistiškai reikšmingai nepranoksta paprastesnių bazinių matų.
KNN laiko eilučių prognozės tyrimuose vis dėlto įsitvirtinęs Euklido atstumas \cite{martinez2019}.
Autoriai pažymi, kad visi jų apžvelgti KNN laiko eilučių prognozavimo darbai naudoja būtent jį, o jų pačių eksperimentuose Euklido ir Manheteno atstumų tikslumas reikšmingai nesiskyrė.

Sąvoka „svertinis KNN“ literatūroje gali reikšti du skirtingus dalykus -- svorius požymiams atstumo funkcijos viduje, kaip (\ref{eq:svertinis-euklidas}) formulėje, ir svorius kaimynams.
Kaimynų svorių variantas literatūroje sutinkamas beveik visur ir dažnai vadinamas tiesiog svertiniu artimiausių kaimynų metodu, nors požymių svorių jame nėra.
Skaitant mokslinius darbus, svarbu tikrinti, kuris variantas turimas omenyje.

%------------------------------------------------------------------

\subsection{Kaimynų skaičius ir prognozės sudarymas}
%------------------------------------------------------------------

Kaimynų skaičius $k$ apibrėžia, į kiek pavyzdžių bus atsižvelgiama priimant sprendimą dėl vertės prognozavimo.
Imant daugiau kaimynų, triukšmas geriau išsividurkina, tačiau tolimesni kaimynai vis mažiau panašūs į dabartinę būseną, todėl prognozė artėja link bendro vidurkio.
Praktikoje $k$ dažniausiai parenkamas iš duomenų.
Išskiriamos trys strategijos: euristinė, kai $k$ prilyginamas kvadratinei šakniai iš mokymo pavyzdžių skaičiaus, $k$ optimizavimas pagal paklaidą atskiroje patikros aibėje ir kelių modelių su skirtingais $k$ prognozių vidurkinimas \cite{martinez2019}.
Eksperimentuose tiksliausias ir kartu greitesnis už optimizavimą buvo kelių modelių derinys.
Elektros kainų prognozėje \cite{troncoso2007} pagal paklaidą optimizuoja tiek lango ilgį, tiek kaimynų skaičių.

Fiksuotas, visiems taškams vienodas $k$ apskritai nėra optimalus, ką teoriškai pagrindžia \cite{anava2016}.
Kiekvieno prognozuojamo taško kaimynystės geometrija skiriasi, todėl autoriai kiekvienam taškui atskirai minimizuoja prognozės paklaidos rėžį pagal kaimynams priskiriamus svorius.
Nustatyta, kad optimalus svoris mažėja tiesiškai didėjant atstumui, o nuo tam tikro atstumo tampa lygus nuliui.
Prognozėje dalyvauja tik dalis kaimynų, ir jų skaičius kinta nuo taško prie taško.

Atrinkus kaimynus, prognozė sudaroma iš jų reikšmių.
Paprasčiausias būdas -- aritmetinis vidurkis, kai visi kaimynai lygiaverčiai:

\begin{equation}
    \hat{y} = \frac{1}{k} \sum_{i=1}^{k} y_i
    \label{eq:kaimynu-vidurkis}
\end{equation}

čia $y_i$ -- $i$-ojo pagal atstumą kaimyno reikšmė.
Regresijoje tai prognozuojamo dydžio reikšmė, o laiko eilutėje -- po kaimyno lango sekanti reikšmė.
Dažnai kaimynams priskiriami svoriai pagal atstumą.
Kuo kaimynas arčiau, tuo didesnė jo įtaka.
Tuomet prognozė yra svertinis kaimynų reikšmių vidurkis:

\begin{equation}
    \hat{y} = \frac{\sum_{i=1}^{k} \omega_i y_i}{\sum_{i=1}^{k} \omega_i}
    \label{eq:svertinis-vidurkis}
\end{equation}

čia $\omega_i \geq 0$ -- $i$-ojo kaimyno svoris.
Svorio funkcijos literatūroje siūlomos įvairios.
\cite{zamani2008} palygina tiesinę, atvirkštinę ir atvirkštinę kvadratinę svorio funkcijas, kuriose svoris skaičiuojamas tiesiogiai iš absoliutaus atstumo.
Tačiau kai artimiausias kaimynas beveik sutampa su prognozuojamu tašku, atvirkštinis atstumas jam suteikia milžinišką svorį ir metodas praktiškai virsta vieno artimiausio kaimyno taisykle.

Šio trūkumo išvengia \cite{dudani976} pasiūlytas svoris, normalizuotas pagal pačią kaimynystę.
Iš $k$ atrinktų kaimynų paimami artimiausio kaimyno atstumas $d_1$ ir tolimiausio kaimyno atstumas $d_k$, o $i$-ojo kaimyno svoris apskaičiuojamas pagal formulę:

\begin{equation}
    \omega_i =
    \begin{cases}
        \dfrac{d_k - d_i}{d_k - d_1}, & d_k \neq d_1 \\[2ex]
        1, & d_k = d_1
    \end{cases}
    \label{eq:wknn-svoris}
\end{equation}

Artimiausias kaimynas gauna svorį, lygų vienetui, tolimiausias -- nuliui, o kitų kaimynų svoriai pasiskirsto visame šiame intervale net tada, kai visų kaimynų atstumai panašūs.
Autorius parodė, kad viename mažos mokymo imties atvejyje svertinė taisyklė esant k>3 klydo rečiau už paprastą balsavimą ir iškėlė prielaidą, jog taip bus ir kitoms mažoms bei vidutinėms imtims \cite{dudani976}.

Dvigubai svertinis variantas \cite{gou2011} normalizuotą svorį papildomai daugina iš kaimyno eilės numerio atvirkštinės reikšmės:

\begin{equation}
    \omega_i =
    \begin{cases}
        \dfrac{d_k - d_i}{d_k - d_1} \cdot \dfrac{1}{i}, & d_k \neq d_1 \\[2ex]
        1, & d_k = d_1
    \end{cases}
    \label{eq:dwknn-svoris}
\end{equation}

Kraštinių kaimynų svoriai lieka tokie patys, tačiau vidutinio atstumo kaimynai nuslopinami, todėl pavieniai netipiški pavyzdžiai kaimynystėje daro mažesnę įtaką.
Autorių eksperimentuose dvylikoje realių duomenų rinkinių šis variantas pasiekė mažiausią vidutinę paklaidą, o tikslumas išliko stabilus plačiame $k$ reikšmių intervale, kai visų trijų palyginamųjų metodų, įskaitant bazinį KNN, tikslumas didėjant $k$ sparčiai krito.
Literatūroje naudojama ir kitokių dvigubo svėrimo variantų, pavyzdžiui, \cite{huang2017} normalizuotą svorį daugina iš atstumais grįsto daugiklio $(d_k + d_1)/(d_k + d_i)$.

Vis dėlto kaimynų svėrimas prognozę gerina ne visada.
\cite{martinez2019} eksperimentuose su 111 laiko eilučių svertinis kaimynų vidurkis buvo statistiškai reikšmingai blogesnis už paprastą aritmetinį vidurkį, todėl galutinėje jų metodikoje kaimynai lygiaverčiai.
Taigi kaimynų svėrimo nauda priklauso nuo duomenų ir turi būti tikrinama.

%------------------------------------------------------------------

\subsection{Atsisakymo mechanizmas}
%------------------------------------------------------------------

Kaimynų atrankos metu taikomas atsisakymo mechanizmas leidžia modeliui nepateikti prognozės, jei nėra pakankamai panašių pavyzdžių.
Modelio, kuris atsisako dalies prognozių, kokybė vertinama dviem rodikliais.
Paklaida tais atvejais, kai prognozė pateikta, ir padengimu -- dalis atvejų, kuriems prognozė apskritai pateikta.
Šie rodikliai konkuruoja, nes griežtesnis atsisakymas mažina paklaidą, bet kartu ir padengimą.

\cite{denis2020} nagrinėja regresiją su atsisakymu.
Optimali atsisakymo taisyklė grindžiama sąlyginės dispersijos slenksčiavimu.
Prognozė pateikiama tik tada, kai sąlyginė dispersija toje srityje neviršija slenksčio, o slenkstis kalibruojamas taip, kad atsisakymų dalis atitiktų iš anksto pasirinktą biudžetą.
Taisyklę autoriai pritaiko būtent KNN prognozavimui ir pažymi, kad atsisakymas regresijoje nagrinėjamas itin retai, nors atmetus dalį nepatikimų atvejų likusiųjų prognozės paklaida sumažėja.
Atsisakymas čia grindžiamas prognozės dispersija, o ne kaimynų panašumu.

Padengimo ir rizikos kompromisas kaip tiesioginis optimizavimo tikslas jau naudojamas giliuosiuose neuroniniuose tinkluose.
\cite{geifman2019} pasiūlytas SelectiveNet mokosi prognozuoti ir atsisakyti vienu metu, minimizuodamas paklaidą su vartotojo nurodytu padengimo apribojimu.
KNN metode padengimas literatūroje naudojamas tik atsisakymo slenksčiui suderinti po to, kai modelis jau sudarytas, o ne renkantis patį $k$ ar panašumo ribą.

%------------------------------------------------------------------

%\subsection{Taikymas laiko eilutėms}
%%------------------------------------------------------------------
%
%Taikant KNN laiko eilutėms, reikia nuspręsti, kaip iš eilutės formuoti pavyzdžius, kaip prognozuoti kelis žingsnius į priekį ir kaip įtraukti kalendorinę informaciją.
%Sistemiškai šiuos klausimus nagrinėja \cite{martinez2019}, pasiūlydami automatinę KNN taikymo laiko eilutėms metodiką.
%Laiko eilutė pati savaime nėra požymių vektorių rinkinys, todėl pavyzdžiai iš jos formuojami langais.
%Slenkant langą per visą istoriją po vieną žingsnį, iš vienos eilutės gaunama daug mokymo pavyzdžių.
%Požymiai yra lango reikšmės, tai yra uždelstos prognozuojamo kintamojo reikšmės, o atsakymu laikoma iškart po lango einanti reikšmė \cite{strategiestimeseries, martinez2019}.
%Toks formavimas remiasi prielaida, kad eilutėje kartojasi šablonai, todėl į dabartinį langą panašių langų tęsiniai parodo tikėtiną tolesnę eigą \cite{martinez2019}.
%Kurias uždelstas reikšmes į langą įtraukti, autoriai parenka nuosekliąja atranka pagal paklaidą patikros aibėje.
%
%Prognozuoti keliems žingsniams į priekį galima trimis strategijomis \cite{strategiestimeseries}.
%Rekursinė strategija prognozuoja po vieną žingsnį, o gautą prognozę naudoja kaip įvestį kitam žingsniui, todėl paklaida gali kauptis.
%Tiesioginė strategija kiekvienam horizonto žingsniui kuria atskirą modelį.
%MIMO strategija (angl. \textit{multiple-input multiple-output}) visą horizontą prognozuoja iš karto -- kaimynų atsakymai yra ne pavienės reikšmės, o ištisi ateities langai.
%\cite{martinez2019} eksperimentuose MIMO buvo reikšmingai netiksliausia, tarp rekursinės ir tiesioginės reikšmingo skirtumo nerasta, todėl dėl paprastumo pasirinkta rekursinė.
%Tuose pačiuose eksperimentuose nustatyta, kad su sezoniškumu KNN susidoroja be išankstinio duomenų apdorojimo, tačiau ilgalaikę tendenciją turinčias eilutes prieš prognozuojant būtina transformuoti, nes prognozė negali išeiti už istorinių reikšmių ribų.
%
%Kalendorinę informaciją į kaimynų paiešką galima įtraukti kaip papildomą požymį.
%Prognozuodami Jungtinės Karalystės elektros paklausą, \cite{alqahtani2013} prie uždelstų reikšmių prideda dvejetainį dienos tipo požymį -- darbo arba ne darbo diena -- ir įtraukia jį į kaimynų paiešką.
%Daugiamatis KNN su šiuo požymiu paklaidą sumažino beveik ketvirtadaliu, palyginti su vienmačiu.
%Išvadose autoriai pažymi, kad įtraukiant daugiau požymių iškiltų jų tarpusavio ryšio problema, kuriai spręsti reikėtų automatinio santykinių požymių svorių nustatymo, o daugiamatėje KNN laiko eilučių regresijoje tai dar turi būti išplėtota \cite{alqahtani2013}.
%Ši pastaba svarbi trečiame skyriuje.
%
%Kaimynų paieškos idėja laiko eilutėse plėtojama ir kitu pavadinimu.
%Meteorologijoje ji vadinama analogų ansambliu \cite{dellemonache2013}: dabartinė skaitmeninio orų modelio prognozė lyginama su istorinėmis to paties modelio prognozėmis, surandamos panašiausios, o tikimybinė prognozė sudaroma iš tuo metu realiai stebėtų reikšmių.
%Panašumas matuojamas per kelis prognostinius kintamuosius, kurių kiekvienam formulėje numatytas svoris, tačiau šiame darbe visi svoriai prilyginti vienetui ir, kaip nurodo patys autoriai, jų optimizuoti nebandyta \cite{dellemonache2013}.
%Šią spragą vėliau užpildė \cite{junk2015}, kurių darbas nagrinėjamas trečiame skyriuje.
%
%Galiausiai KNN laiko eilutėse jungiamas su signalo apdorojimo metodais.
%\cite{sudheer2015} trumpalaikei elektros apkrovos prognozei eilutę bangelių transformacija išskaido į glodžiąją ir svyravimų dalis: glodžioji prognozuojama Holto ir Vinterso eksponentinio glodinimo metodu, o svyravimų dalis -- svertiniu artimiausių kaimynų metodu.
%Tokie hibridai rodo, kad KNN vieta gali būti ir šalia kitų metodų, apdorojant tą signalo dalį, kurioje panašumo paieška veikia geriausiai.
%
%Apibendrinant, KNN veikimą apibrėžia keli tarpusavyje susiję sprendimai: atstumo funkcija, kaimynų skaičius, kaimynų svoriai ir pavyzdžių formavimas.
%Kaimynų svoriai literatūroje ištirti išsamiai, o požymių svoriai atstumo funkcijoje, nuo kurių tiesiogiai priklauso pati kaimynų atranka, laiko eilučių prognozėje nagrinėjami retai.
%Kitas skyrius skirtas būtent jiems.
